% SPDX-License-Identifier: CC-BY-SA-4.0
% Copyright (c) 2026 Rafael Beraldo
% Parte 2: o básico do básico (RUG)
% Incluído por presentation.tex; compile a partir dele.

% =============================================================================
\section{O básico do básico}
% =============================================================================

\begin{frame}{Uma solução simples: a Radiocomunicação de Uso Geral (RUG)}
  \begin{columns}[T,onlytextwidth]
    \column{0.55\textwidth}
      Para comunicação na trilha ou na escalada sem licença, use um \alert{kit de rádios de uso geral} (RUG).

      \vspace{.5em}
      \begin{uncoverenv}<2->
        \textbf{Prós}
        \begin{itemize}
        \item Kits relativamente baratos (R\$ 200–400)
        \item Legalizado \citep{anatel2017}
        \item Resolve o problema
        \end{itemize}
      \end{uncoverenv}

      \begin{uncoverenv}<3->
        \textbf{Contras}
        \begin{itemize}
        \item Baixa potência (500\,mW)
        \item Não há repetidoras
        \end{itemize}
      \end{uncoverenv}

    \column{0.4\textwidth}
      % Sem especificação de overlay: visível desde o passo 1
      \foto[width=\linewidth, height=5.5cm, keepaspectratio]{talkabout}{Kit do Talkabout T110br}
  \end{columns}
\end{frame}

\begin{frame}[standout]
  Muito obrigado! Perguntas? \\[.5em] \pause Mas, se quisermos ir além do básico, precisamos nos aprofundar no \alert{radioamadorismo}.
\end{frame}

%%% Local Variables:
%%% TeX-master: "../presentation"
%%% fill-column: 100
%%% End:
